\documentclass[11pt,a4paper]{article}

\usepackage[utf8]{inputenc}
\usepackage[T1]{fontenc}
\usepackage[swedish]{babel}
\usepackage{lmodern}
\usepackage[margin=2.2cm]{geometry}
\usepackage{xcolor}
\usepackage[most]{tcolorbox}
\usepackage{multicol}
\usepackage{enumitem}
\usepackage{microtype}

\definecolor{ink}{HTML}{20313A}
\definecolor{teal}{HTML}{167C80}
\definecolor{mist}{HTML}{EFF6F4}
\definecolor{line}{HTML}{D5E4E1}

\pagestyle{empty}
\setlength{\parindent}{0pt}
\setlength{\parskip}{5pt}
\setlist[itemize]{leftmargin=1.4em,itemsep=2pt,topsep=3pt}

\newcommand{\tagname}[1]{\texttt{\textless\detokenize{#1}\textgreater}}
\newcommand{\cssname}[1]{\texttt{\detokenize{#1}}}

\begin{document}

{\color{teal}\rule{\linewidth}{2pt}}
\vspace{0.35cm}
{\color{ink}\Huge\bfseries HTML och CSS}\par
{\large Sammanställning av genomgångarna i HT26}\par
\vspace{0.2cm}
{\small Unika taggar, CSS-egenskaper och pseudo-selektorer som förekommer i materialet.}
\vspace{0.4cm}

\begin{tcolorbox}[
  colback=mist,
  colframe=line,
  boxrule=0.6pt,
  arc=2mm,
  left=4mm,right=4mm,top=3mm,bottom=3mm
]
\section*{HTML-taggar}
\begin{multicols}{3}
\begin{itemize}
  \item \tagname{a}
  \item \tagname{aside}
  \item \tagname{body}
  \item \tagname{br}
  \item \tagname{button}
  \item \tagname{code}
  \item \tagname{datalist}
  \item \tagname{details}
  \item \tagname{div}
  \item \tagname{fieldset}
  \item \tagname{footer}
  \item \tagname{form}
  \item \tagname{h1}
  \item \tagname{h2}
  \item \tagname{h3}
  \item \tagname{h6}
  \item \tagname{head}
  \item \tagname{header}
  \item \tagname{html}
  \item \tagname{iframe}
  \item \tagname{img}
  \item \tagname{input}
  \item \tagname{label}
  \item \tagname{legend}
  \item \tagname{li}
  \item \tagname{link}
  \item \tagname{main}
  \item \tagname{meta}
  \item \tagname{meter}
  \item \tagname{nav}
  \item \tagname{ol}
  \item \tagname{option}
  \item \tagname{p}
  \item \tagname{progress}
  \item \tagname{section}
  \item \tagname{select}
  \item \tagname{span}
  \item \tagname{style}
  \item \tagname{summary}
  \item \tagname{table}
  \item \tagname{td}
  \item \tagname{textarea}
  \item \tagname{th}
  \item \tagname{title}
  \item \tagname{tr}
  \item \tagname{ul}
\end{itemize}
\end{multicols}
\end{tcolorbox}

\begin{tcolorbox}[
  colback=white,
  colframe=line,
  boxrule=0.6pt,
  arc=2mm,
  left=4mm,right=4mm,top=3mm,bottom=3mm
]
\section*{CSS-egenskaper}
\begin{multicols}{3}
\begin{itemize}
  \item \cssname{align-items}
  \item \cssname{background}
  \item \cssname{background-color}
  \item \cssname{border}
  \item \cssname{border-radius}
  \item \cssname{bottom}
  \item \cssname{box-shadow}
  \item \cssname{box-sizing}
  \item \cssname{color}
  \item \cssname{cursor}
  \item \cssname{display}
  \item \cssname{flex-direction}
  \item \cssname{flex-grow}
  \item \cssname{flex-wrap}
  \item \cssname{font-family}
  \item \cssname{font-size}
  \item \cssname{font-weight}
  \item \cssname{gap}
  \item \cssname{grid-area}
  \item \cssname{grid-column}
  \item \cssname{grid-row}
  \item \cssname{grid-template-areas}
  \item \cssname{grid-template-columns}
  \item \cssname{grid-template-rows}
  \item \cssname{height}
  \item \cssname{justify-content}
  \item \cssname{left}
  \item \cssname{list-style}
  \item \cssname{list-style-type}
  \item \cssname{margin}
  \item \cssname{margin-bottom}
  \item \cssname{min-height}
  \item \cssname{min-width}
  \item \cssname{overflow}
  \item \cssname{padding}
  \item \cssname{position}
  \item \cssname{right}
  \item \cssname{text-align}
  \item \cssname{text-decoration}
  \item \cssname{top}
  \item \cssname{transform}
  \item \cssname{transition}
  \item \cssname{width}
\end{itemize}
\end{multicols}

\subsection*{Egna CSS-variabler}
\cssname{--accent-color}, \cssname{--page-background},
\cssname{--primary-color}, \cssname{--secondary-color},
\cssname{--size}, \cssname{--text-color}
\end{tcolorbox}

\begin{tcolorbox}[
  colback=mist,
  colframe=line,
  boxrule=0.6pt,
  arc=2mm,
  left=4mm,right=4mm,top=3mm,bottom=3mm
]
\section*{Pseudo-selektorer}
\begin{itemize}
  \item \cssname{:hover} --- pseudo-klass som används vid hovring.
  \item \cssname{:root} --- pseudo-klass som väljer dokumentets rotelement.
\end{itemize}
Inga pseudo-element med dubbelkolon, till exempel \cssname{::before}, hittades i materialet.
\end{tcolorbox}

\vfill
{\color{teal}\rule{\linewidth}{0.8pt}}\\[-2pt]
{\small Sammanställd från HTML- och CSS-filerna i HT26.}

\end{document}